%%
%% This is file `mathfont_example_roboto.tex',
%% generated with the docstrip utility.
%%
%% The original source files were:
%%
%% mathfont_code.dtx  (with options: `roboto')
%% 
%% This file is from version 2.3 of the free and open-source
%% LaTeX package "mathfont," released September 2023, to be used
%% with the XeTeX or LuaTeX engines. (As of version 2.0, LuaTeX
%% is recommended.)
%% 
%% Copyright 2018-2023 by Conrad Kosowsky
%% 
%% This Work may be used, distributed, and modified under the
%% terms of the LaTeX Public Project License, version 1.3c or
%% any later version. The most recent version of this license
%% is available online at
%% 
%%           https://www.latex-project.org/lppl/.
%% 
%% This Work has the LPPL status "maintained," and the current
%% maintainer is the package author, Conrad Kosowsky. He can
%% be reached at kosowsky.latex@gmail.com.
%% 
%% PLEASE KNOW THAT THIS FREE SOFTWARE IS PROVIDED WITHOUT
%% ANY WARRANTY. SPECIFICALLY, THE "NO WARRANTY" SECTION OF
%% THE LATEX PROJECT PUBLIC LICENSE STATES THE FOLLOWING:
%% 
%% THERE IS NO WARRANTY FOR THE WORK. EXCEPT WHEN OTHERWISE
%% STATED IN WRITING, THE COPYRIGHT HOLDER PROVIDES THE WORK
%% `AS IS’, WITHOUT WARRANTY OF ANY KIND, EITHER EXPRESSED
%% OR IMPLIED, INCLUDING, BUT NOT LIMITED TO, THE IMPLIED
%% WARRANTIES OF MERCHANTABILITY AND FITNESS FOR A PARTICULAR
%% PURPOSE. THE ENTIRE RISK AS TO THE QUALITY AND PERFORMANCE
%% OF THE WORK IS WITH YOU. SHOULD THE WORK PROVE DEFECTIVE,
%% YOU ASSUME THE COST OF ALL NECESSARY SERVICING, REPAIR, OR
%% CORRECTION.
%% 
%% IN NO EVENT UNLESS REQUIRED BY APPLICABLE LAW OR AGREED
%% TO IN WRITING WILL THE COPYRIGHT HOLDER, OR ANY AUTHOR
%% NAMED IN THE COMPONENTS OF THE WORK, OR ANY OTHER PARTY
%% WHO MAY DISTRIBUTE AND/OR MODIFY THE WORK AS PERMITTED
%% ABOVE, BE LIABLE TO YOU FOR DAMAGES, INCLUDING ANY GENERAL,
%% SPECIAL, INCIDENTAL OR CONSEQUENTIAL DAMAGES ARISING OUT
%% OF ANY USE OF THE WORK OR OUT OF INABILITY TO USE THE WORK
%% (INCLUDING, BUT NOT LIMITED TO, LOSS OF DATA, DATA BEING
%% RENDERED INACCURATE, OR LOSSES SUSTAINED BY ANYONE AS A
%% RESULT OF ANY FAILURE OF THE WORK TO OPERATE WITH ANY
%% OTHER PROGRAMS), EVEN IF THE COPYRIGHT HOLDER OR SAID
%% AUTHOR OR SAID OTHER PARTY HAS BEEN ADVISED OF THE
%% POSSIBILITY OF SUCH DAMAGES.
%% 
%% For more information, see the LaTeX Project Public License.
%% Derivative works based on this package may come with their
%% own license or terms of use, and the package author is not
%% responsible for any third-party software.
%% 
%% For more information, see mathfont_code.dtx. Happy TeXing!
%% 
\ifx\directlua\undefined
  \PackageError{mathfont}
  {\MessageBreak
  LuaLaTeX is recommended for\MessageBreak
  mathfont_example_roboto.tex}
  {It's recommended that you typset this file\MessageBreak
  with LuaTeX. You can use XeLaTeX if you want,\MessageBreak
  but things will turn out weird. To resolve\MessageBreak
  this error, typeset the file with LuaLaTeX.}
\fi
\documentclass[12pt]{article}
\usepackage[margin=1in]{geometry}
\usepackage{innerscript}
\usepackage{multicol}
\usepackage{amsmath}
\DeclareMathOperator{\Res}{Res}
\usepackage{mathfont}
\mathfont[radical]{Overpass}
\setfont{Roboto}
\mathfont[]{STIXGeneral}
\DeclareSymbolFont{Mupright2}{TU}{STIXGeneral}{m}{n}
\DeclareMathSymbol{\otimes}{\mathbin}{Mupright2}{"2297}
\RuleThicknessFactor{1900}
\IntegralItalicFactor{500}
\SurdHorizontalFactor{800}
\parindent=0pt\relax
\begin{document}

\def\footnote#1{}
\def\showabstract{0}
\let\textsf\relax
\let\ttfamily\relax
\def\documentname{Example---Roboto}

\input mathfont_heading.tex

This is Roboto with Overpass for the radical signs and STIXGeneral for the tensor product. ``Testing. Testing.'' Brown foxes quickly jump over dazzling does and harts. This document shows an example of mathfont in action. Unfortunately, there are many more equations in the world than I have space for here. Nevertheless, I hope I hit some of the highlights. Happy \TeX ing!

%%
%% This is file `mathfont_equations.tex',
%% generated with the docstrip utility.
%%
%% The original source files were:
%%
%% mathfont_code.dtx  (with options: `equations')
%% 
%% This file is from version 2.3 of the free and open-source
%% LaTeX package "mathfont," released September 2023, to be used
%% with the XeTeX or LuaTeX engines. (As of version 2.0, LuaTeX
%% is recommended.)
%% 
%% Copyright 2018-2023 by Conrad Kosowsky
%% 
%% This Work may be used, distributed, and modified under the
%% terms of the LaTeX Public Project License, version 1.3c or
%% any later version. The most recent version of this license
%% is available online at
%% 
%%           https://www.latex-project.org/lppl/.
%% 
%% This Work has the LPPL status "maintained," and the current
%% maintainer is the package author, Conrad Kosowsky. He can
%% be reached at kosowsky.latex@gmail.com.
%% 
%% PLEASE KNOW THAT THIS FREE SOFTWARE IS PROVIDED WITHOUT
%% ANY WARRANTY. SPECIFICALLY, THE "NO WARRANTY" SECTION OF
%% THE LATEX PROJECT PUBLIC LICENSE STATES THE FOLLOWING:
%% 
%% THERE IS NO WARRANTY FOR THE WORK. EXCEPT WHEN OTHERWISE
%% STATED IN WRITING, THE COPYRIGHT HOLDER PROVIDES THE WORK
%% `AS IS’, WITHOUT WARRANTY OF ANY KIND, EITHER EXPRESSED
%% OR IMPLIED, INCLUDING, BUT NOT LIMITED TO, THE IMPLIED
%% WARRANTIES OF MERCHANTABILITY AND FITNESS FOR A PARTICULAR
%% PURPOSE. THE ENTIRE RISK AS TO THE QUALITY AND PERFORMANCE
%% OF THE WORK IS WITH YOU. SHOULD THE WORK PROVE DEFECTIVE,
%% YOU ASSUME THE COST OF ALL NECESSARY SERVICING, REPAIR, OR
%% CORRECTION.
%% 
%% IN NO EVENT UNLESS REQUIRED BY APPLICABLE LAW OR AGREED
%% TO IN WRITING WILL THE COPYRIGHT HOLDER, OR ANY AUTHOR
%% NAMED IN THE COMPONENTS OF THE WORK, OR ANY OTHER PARTY
%% WHO MAY DISTRIBUTE AND/OR MODIFY THE WORK AS PERMITTED
%% ABOVE, BE LIABLE TO YOU FOR DAMAGES, INCLUDING ANY GENERAL,
%% SPECIAL, INCIDENTAL OR CONSEQUENTIAL DAMAGES ARISING OUT
%% OF ANY USE OF THE WORK OR OUT OF INABILITY TO USE THE WORK
%% (INCLUDING, BUT NOT LIMITED TO, LOSS OF DATA, DATA BEING
%% RENDERED INACCURATE, OR LOSSES SUSTAINED BY ANYONE AS A
%% RESULT OF ANY FAILURE OF THE WORK TO OPERATE WITH ANY
%% OTHER PROGRAMS), EVEN IF THE COPYRIGHT HOLDER OR SAID
%% AUTHOR OR SAID OTHER PARTY HAS BEEN ADVISED OF THE
%% POSSIBILITY OF SUCH DAMAGES.
%% 
%% For more information, see the LaTeX Project Public License.
%% Derivative works based on this package may come with their
%% own license or terms of use, and the package author is not
%% responsible for any third-party software.
%% 
%% For more information, see mathfont_code.dtx. Happy TeXing!
%% 
\multicolsep=0pt\relax
\bigskip
\begin{multicols*}{2}

Black-Scholes Equation
\[
\frac{\partial V}{\partial t}+
  \frac12\sigma^2S^2\frac{\partial^2V}{\partial S^2}
  =rV-rS\frac{\partial V}{\partial X}
\]

\vfill

Cardano's Formula/Cubic Formula
\begin{align*}
t_i&=\omega_i\sqrt[3]{-\frac q2+\sqrt{\frac{q^2}4+\frac{p^3}{27}}}\\
  &\qquad\qquad{}+\omega_i^2\sqrt[3]{-\frac q2-\sqrt{\frac{q^2}4+\frac{p^3}{27}}}
\end{align*}

\vfill

Einstein's Field Equation (General Relativity)
\[
R_{\mu\nu}-\frac12Rg_{\mu\nu}+\Lambda g_{\mu\nu}=\frac{8\pi G}{c^4}T_{\mu\nu}
\]

\vfill

First Isomorphism Theorem
\[
\phi(X)\cong X/\ker(\phi)
\]

\vfill

Gauss-Bonnet Formula
\[
\int_MK\ dA+\int_{\partial M}k_g\ ds=2\pi\chi(M)
\]

\vfill

Maxwell's Equations
\begin{align*}
\nabla\bullet\mathbf E&=\frac\rho{\epsilon_0}&
\nabla\bullet\mathbf B&=0\\
\nabla\times\mathbf E&=-\frac{\partial\mathbf B}{\partial t}&
\nabla\times\mathbf B&=\mu_0\left(\mathbf J+
  \epsilon_0\frac{\partial\mathbf E}{\partial t}\right)
\end{align*}

\vfill

Michaelis-Menten Model
\bgroup
\belowdisplayskip=0pt\relax
\belowdisplayshortskip=0pt\relax
\[
v=\frac{d[P]}{dt}=V\frac{[S]}{K_M+[S]}
\]
\par\egroup
\hrule height 0pt
\columnbreak
%% next column begins here

Navier-Stokes Equation
\begin{align*}
\frac{\partial}{\partial t}(\rho\mathbf{u})+\nabla\bullet
  (\rho \mathbf{u}\otimes \mathbf{u})={}
  &{-}\nabla\bar p+\mu\nabla^2\mathbf{u}\\
  &+\frac13\mu\nabla(\nabla\bullet u)+
  \rho\mathbf{g}
\end{align*}

\vfill

Quadratic Formula
\[
x=\frac{-b\pm\sqrt{b^2-4ac}}{2a}
\]

\vfill

Ramanujan's Approximation for $\Gamma$
\[
\Gamma(1+x)\approx\sqrt\pi\,x^xe^{-x}\,\sqrt[6]{8x^3+4x^2+x+\frac1{30}}
\]

\vfill

Residue Theorem
\[
\frac1{2i\pi}\int_{\gamma}f(z)\ dz=\sum_{k=1}^n\Res_{a_k}(f)
\]

\vfill

Riemann Zeta Function
\begin{align*}
\zeta(z)&=\sum_{i=1}^\infty\frac1{z^i}
  =\frac1{\Gamma(z)}\int_0^\infty\frac{x^{z-1}}{e^x-1}\ dx\\
  &=2^z\pi^{z-1}\sin\left(\frac{\pi z}2\right)\,\Gamma(1-z)\,\zeta(1-z)
\end{align*}

\vfill

Schrodinger Equation
\[
i\hbar\frac d{dt}|\Psi(t)\fakerangle=\hat H|\Psi(t)\fakerangle
\]

\vfill

Lorentz Transformation (Special Relativity)
\belowdisplayskip=0pt\relax
\belowdisplayshortskip=0pt\relax
\[
t'=\left(t-\frac{vx}{c^2}\right)\frac1{\sqrt{1-\frac{v^2}{c^2}}}
\]
\hrule height 0pt

\end{multicols*}
\endinput
%%
%% End of file `mathfont_equations.tex'.


\end{document}
\endinput
%%
%% End of file `mathfont_example_roboto.tex'.
